% Vietnamese translation for OI Public Library
% Source-File: IOI中国国家候选队论文/2008/Day1/5.周梦宇《码之道——浅谈信息学竞赛中的编码与译码问题》/码之道.pptx
\documentclass[11pt]{article}
\usepackage{oipl-vn}

\newcommand{\Oh}{\mathcal{O}}
\newcommand{\permrank}{\operatorname{rank}}

\title{Đạo của mã\\{\large Bàn về mã hóa và giải mã trong thi tin học}}
\author{Zhou Mengyu (Moony Chou)\\Trường Trung học số 7 Thành Đô}
\date{}

\begin{document}
\maketitle

\origtitle{The Tao of Coding}
\sourcepath{IOI 2008 / Zhou Mengyu / The Tao of Coding}

\section{Thông tin và mã hóa}

Thông tin hiện diện ở khắp nơi dưới dạng các bit nhỏ bé \(0\) và \(1\).
Bài báo kinh điển \emph{The Mathematical Theory of Communication} của
Claude E. Shannon đặt nền móng cho lý thuyết thông tin. Điện báo, điện
thoại, phát thanh, điều khiển từ xa, radar và nhiều hệ thống khác đều có
thể được nhìn như các hệ truyền thông.

Một mô hình truyền thông cơ bản gồm:
\[
  \text{nguồn tin}
  \longrightarrow
  \text{bộ mã hóa}
  \longrightarrow
  \text{kênh}
  \longrightarrow
  \text{bộ giải mã}
  \longrightarrow
  \text{đích nhận}.
\]
Trong quá trình truyền, nguồn nhiễu có thể tác động lên tín hiệu trong kênh.
Vì thế ta thường tách các nhiệm vụ thành mã hóa nguồn, mã hóa mật mã, mã hóa
kênh, rồi ở phía nhận thực hiện giải mã kênh, giải mật mã và giải mã nguồn.

Trong ngữ cảnh thi tin học, ``mã hóa'' thường là biến một đối tượng tổ hợp
thành một số thứ tự hoặc một biểu diễn ngắn gọn; ``giải mã'' là dựng lại đối
tượng từ số thứ tự hoặc biểu diễn ấy.

\section{Ví dụ: thứ tự từ điển của hoán vị}

Cho một tập ký tự \(S\), với giả thiết \(\lvert S\rvert! < 2^{63}\).
Xét hai bài toán:
\begin{enumerate}
  \item Tìm hoán vị đứng thứ \(k\) theo thứ tự từ điển trong tất cả các hoán
  vị của \(S\).
  \item Cho một hoán vị của \(S\), hỏi nó đứng thứ mấy theo thứ tự từ điển.
\end{enumerate}

Cách trực tiếp là sắp xếp các phần tử của \(S\) để lấy hoán vị nhỏ nhất theo
thứ tự từ điển, rồi liên tục sinh hoán vị kế tiếp bằng thuật toán
\texttt{next\_permutation}. Với bài toán giải mã, gọi thuật toán \(k-1\) lần
sẽ thu được chuỗi cần tìm. Với bài toán mã hóa, ta phải sinh liên tiếp và so
sánh cho đến khi gặp chuỗi đã cho.

Độ phức tạp của cách này là \(\Oh(\lvert S\rvert k)\), trong khi \(k\) có
thể lớn đến \(\lvert S\rvert!\).

\subsection{\texorpdfstring{\texttt{next\_permutation}}{next permutation}}

Thuật toán sinh hoán vị kế tiếp có thể viết như sau:
\begin{enumerate}
  \item Với chuỗi \(s\), tìm từ phải sang trái chỉ số đầu tiên \(i\) thỏa
  \(s_i < s_{i+1}\).
  \item Đổi \(s_i\) với phần tử nhỏ nhất ở bên phải nó nhưng vẫn lớn hơn
  \(s_i\).
  \item Đảo ngược hậu tố đứng sau vị trí \(i\).
\end{enumerate}
Mỗi lần sinh mất \(\Oh(\lvert s\rvert)\). Thuật toán này rất tiện dụng, nhưng
không đủ khi ta cần nhảy thẳng đến một thứ hạng rất lớn.

\section{Tư tưởng đếm theo lớp}

Để mã hóa một chuỗi, tức tìm thứ hạng của nó, ta đang hỏi: ``có bao nhiêu đối
tượng không lớn hơn \(x\)?'' Đây là một bài toán đếm. Thứ tự từ điển được
quyết định bởi vị trí đầu tiên mà hai chuỗi khác nhau, nên ta có thể phân lớp
theo từng ký tự từ trái sang phải và cộng dồn số phần tử của các lớp nhỏ hơn.

Ví dụ với hoán vị \(42153\) của tập \(\{1,2,3,4,5\}\):
\[
\begin{array}{c|c|c}
\text{tiền tố nhỏ hơn} & \text{số lựa chọn} & \text{số hoán vị} \\
\hline
1????,\ 2????,\ 3???? & 3 & 3\cdot 4! = 72 \\
41??? & 1 & 1\cdot 3! = 6 \\
42135 & 2 & 2\cdot 0! = 2
\end{array}
\]
Vì vậy có \(72+6+2=80\) hoán vị đứng trước \(42153\) nếu đánh số từ \(0\).
Nếu đánh số từ \(1\), thứ hạng là \(81\).

Quá trình giải mã thứ hạng \(80\) làm ngược lại. Trước hết mỗi lớp
\(1????,2????,3????,4????\) có \(4!=24\) phần tử. Tổng dồn lần lượt là
\(24,48,72,96\), nên tiền tố đầu tiên là \(4\). Sau đó xét các lớp
\(41???,42???,\ldots\); mỗi lớp còn \(3!=6\) phần tử, và ta tiếp tục cho đến
khi dựng xong toàn bộ chuỗi.

Nếu độ dài chuỗi cần tìm là \(n\), kích thước tập ký tự là \(m\), thì mỗi vị
trí có thể phải cộng dồn qua tối đa \(m\) ký tự. Nếu thao tác đếm cho từng
lớp mất \(\Oh(m)\), tổng thời gian là
\[
  \Oh(nm^2).
\]
Con số này tốt hơn rất nhiều so với duyệt tuần tự \(\Oh(n\cdot n!)\).

\section{Các công cụ thường gặp}

Các bài toán mã hóa và giải mã trong OI thường dùng phối hợp nhiều nhóm ý
tưởng:
\begin{itemize}
  \item số học và đếm tổ hợp, ví dụ \emph{Twofive} trong IOI 2001 hoặc
  \emph{Hexadecimal Numbers} trong CEOI 1997;
  \item quy hoạch động truy hồi, liệt kê, tìm kiếm, sắp xếp, tham lam và xây
  dựng, ví dụ \emph{CUDAK}, \emph{Zip}, hoặc \emph{A decorative fence};
  \item các cấu trúc dữ liệu đặc biệt như heap và cây Fenwick.
\end{itemize}

Một số ứng dụng tiêu biểu là mã Huffman, mã Gray, mã Prüfer và hàm băm.

\section{\texorpdfstring{Mã Prüfer}{Mã Prüfer}}

Một \term{cây gán nhãn} là cây có \(n\) đỉnh được gán nhãn từ \(1\) đến
\(n\). Định lý Cayley nói rằng số cây gán nhãn khác nhau trên \(n\) đỉnh là
\[
  n^{n-2}.
\]
Năm 1918, Heinz Prüfer đưa ra một chứng minh xây dựng khác thông qua mã
Prüfer: một song ánh giữa cây gán nhãn \(n\) đỉnh và các dãy số nguyên độ dài
\(n-2\), trong đó mỗi phần tử thuộc đoạn \([1,n]\).

\subsection{Mã hóa}

Với cây gán nhãn \(T\) có \(n\) đỉnh, lặp lại thao tác sau cho đến khi cây
chỉ còn hai đỉnh:
\begin{enumerate}
  \item chọn lá có nhãn nhỏ nhất trong cây hiện tại;
  \item ghi lại nhãn của đỉnh kề với lá đó;
  \item xóa lá khỏi cây.
\end{enumerate}
Nhãn ghi được ở lần xóa thứ \(i\) chính là phần tử thứ \(i\) của dãy Prüfer.

Ví dụ trên slide cho một cây có \(8\) đỉnh và dãy Prüfer
\[
  (2,1,3,3,5,7).
\]
Các ngoặc trên hình gốc đánh dấu tuần tự những lá bị xóa trong quá trình mã
hóa.

\subsection{Giải mã}

Cho dãy Prüfer
\[
  (a_1,a_2,\ldots,a_{n-2}).
\]
Ta thêm một phần tử phụ \(n\) ở cuối để dễ mô tả, rồi tại bước \(i\), chọn
nhãn nhỏ nhất chưa bị xóa và không xuất hiện trong phần còn lại
\[
  a_i,a_{i+1},\ldots,a_{n-2}.
\]
Gọi nhãn đó là \(b_i\). Nối cạnh \((b_i,a_i)\), rồi đánh dấu \(b_i\) là đã
bị xóa. Sau \(n-2\) bước, hai nhãn còn lại được nối với nhau. Cách làm này
khôi phục đúng cây ban đầu.

\subsection{Ứng dụng đếm cây}

Bài \emph{Tree Counting} trong HNOI 2004 Day 2 có dạng: cho một cây trên
\(n\) đỉnh \(v_1,v_2,\ldots,v_n\), biết bậc của \(v_i\) là \(d_i\). Hỏi có
bao nhiêu cây khác nhau thỏa điều kiện này?

Trong mã Prüfer của một cây gán nhãn, nhãn \(i\) xuất hiện đúng
\(d_i-1\) lần. Do đó số cây gán nhãn \(n\) đỉnh với dãy bậc đã biết là số
cách sắp xếp đa tập gồm \(d_i-1\) bản sao của nhãn \(i\):
\[
  \frac{(n-2)!}{\prod_{i=1}^{n}(d_i-1)!},
\]
với điều kiện \(\sum_{i=1}^{n} d_i = 2n-2\) và mọi \(d_i\ge 1\).

\section{Một vài ứng dụng khác của mã hóa}

Mã hóa và giải mã xuất hiện trong rất nhiều hệ thống thực tế:
\begin{itemize}
  \item hệ mật mã khóa công khai RSA;
  \item chuẩn nén ảnh JPEG;
  \item MD5 và SHA;
  \item ISBN, mã bưu chính và số điện thoại;
  \item mã kiểm tra vòng CRC;
  \item mã hóa độ dài loạt, mã MH và thuật toán nén LZW;
  \item truyền thông mạng, liên lạc an toàn và kiểm soát lỗi.
\end{itemize}

\section{Kết luận}

Mã hóa và giải mã là một nguồn ý tưởng quan trọng trong xử lý thông tin.
Trong thi tin học, chúng thường biến bài toán thành một câu hỏi đếm, một
song ánh tổ hợp, hoặc một cách biểu diễn giúp chuyển đổi giữa đối tượng và
số thứ tự. Khi gặp một bài toán có ``thứ hạng'', ``khôi phục'', ``mã'', hoặc
``biểu diễn'', nên nghĩ đến việc phân lớp đối tượng, đếm lớp và xây dựng song
ánh thích hợp.

Tác giả kết thúc bài trình bày bằng lời cảm ơn gia đình, thầy Zhang Junliang,
các huấn luyện viên, bạn bè và CCF vì cơ hội được trình bày.

\end{document}
